After the chamber configurations are determined, as a second step we implement a local search phase to fine-tune the chamber configuration decisions. We use the scheduling procedure to form a schedule according to the given chamber configurations and the sample-size vector. The normalized workload of chamber \(c\) is \(\lambda_c=B_c/K_c\), where \(B_c\) is the total duration of the tests assigned to that chamber and \(K_c\) is its number of stations. A chamber is \emph{highly loaded} if \(\lambda_c\) is large and \emph{lightly loaded} if \(\lambda_c\) is small. We consider exchanging the configurations of a highly loaded chamber and a lightly loaded chamber with \(\lambda_c>\lambda_{c'}\) (without violating \eqref{eq:assign_one}--\eqref{eq:assign_hum}). For each feasible swap we call the scheduling procedure and update the assignment if the solution improves. If no swap improves the solution, we reassign a lightly loaded chamber to the configuration of a highly loaded chamber. If that move improves the solution, we update it. These two moves are implemented sequentially as long as the solution improves, or until \(B_{\mathrm{room}}\) calls of the scheduling procedure have been used. The steps are provided in Algorithm~\ref{alg:chamber_assignment}.

\begin{algorithm}[H]
\caption{\textsc{ChamberAssignment}()}
\label{alg:chamber_assignment}
\begin{algorithmic}[1]
\State \textbf{Input:} sample-size vector \(\mathbf{s}\)
\State \textbf{Output:} chamber configuration assignment \(\mathbf{x}\)
\State \(\mathbf{x}\leftarrow\Call{MIP-CCD}{\mathbf{s}}\) \label{step:ca-exact}
\State \((\Pi,T^{tot})\leftarrow\Call{Schedule}{\mathbf{x},\mathbf{s}}\) \label{step:ca-eval}
\Repeat
    \State For pairs \((c,c')\) with \(\lambda_c>\lambda_{c'}\), set \(\mathbf{x}_{cand}\leftarrow\Call{Swap}{\mathbf{x},c,c'}\); if \(\mathbf{x}_{cand}\) is feasible, \((\Pi_{cand},T^{tot}_{cand})\leftarrow\Call{Schedule}{\mathbf{x}_{cand},\mathbf{s}}\) \label{step:ca-swap}
    \If{\((T^{tot},\mathbf{1}^{\mathsf{T}}\mathbf{s})\succ_{\mathrm{lex}}(T^{tot}_{cand},\mathbf{1}^{\mathsf{T}}\mathbf{s})\)} \label{step:ca-accept}
        \State \((\mathbf{x},\Pi,T^{tot})\leftarrow(\mathbf{x}_{cand},\Pi_{cand},T^{tot}_{cand})\)
    \EndIf
    \If{no swap was accepted}
        \State For pairs \((c,c')\) with \(\lambda_c>\lambda_{c'}\), set \(\mathbf{x}_{cand}\leftarrow\Call{Move}{\mathbf{x},c,c'}\); if \(\mathbf{x}_{cand}\) is feasible, \((\Pi_{cand},T^{tot}_{cand})\leftarrow\Call{Schedule}{\mathbf{x}_{cand},\mathbf{s}}\) \label{step:ca-move}
        \If{\((T^{tot},\mathbf{1}^{\mathsf{T}}\mathbf{s})\succ_{\mathrm{lex}}(T^{tot}_{cand},\mathbf{1}^{\mathsf{T}}\mathbf{s})\)}
            \State \((\mathbf{x},\Pi,T^{tot})\leftarrow(\mathbf{x}_{cand},\Pi_{cand},T^{tot}_{cand})\)
        \EndIf
    \EndIf
\Until{no improvement or \(B_{\mathrm{room}}\) calls of \(\textsc{Schedule}()\)} \label{step:ca-stop}
\State \Return \(\mathbf{x}\)
\end{algorithmic}
\end{algorithm}

In Algorithm~\ref{alg:chamber_assignment}, MIP-CCD(\(\mathbf{s}\)) gives the initial assignment and \(\textsc{Schedule}()\) evaluates it (Steps~\ref{step:ca-exact}--\ref{step:ca-eval}). We then try \(\textsc{Swap}\) over pairs with \(\lambda_c>\lambda_{c'}\) (Step~\ref{step:ca-swap}) and accept a feasible neighbor if it improves the lexicographic objective (Step~\ref{step:ca-accept}). If no swap is accepted, we try \(\textsc{Move}\) in the same way (Step~\ref{step:ca-move}). The search stops when a pass yields no improvement or \(B_{\mathrm{room}}\) calls of \(\textsc{Schedule}()\) have been used (Step~\ref{step:ca-stop}).
